\clearpage
\fancyhead[L]{\small INSTRUMENTS.}
\hypertarget{native-427-part-lower}{}
\begin{center}
\Large INSTRUMENTS.
\end{center}



\begin{center}
Göttingische gelehrte Anzeigen. No. 167. pp. 1665..1669. 18 October 1819.
\end{center}





In issue 127 of this journal last year [p. 410 of this volume], we gave a description of \textsc{Repsold}’s meridian circle and some of the first observations made with it at this observatory. We now give a brief account of \textsc{Reichenbach}’s \emph{transit telescope}, the second principal instrument, whose installation was completed recently. Since the design of transit telescopes in use at the current state of observational practice is always the same in all essential respects, we may confine ourselves here to the features peculiar to ours. The telescope has a focal length of six Paris feet and an aperture of 52 lines. In aperture it is surpassed only by the new transit telescope at Greenwich Observatory, which at a focal length of 10 feet has an aperture of 5 English inches. The axis measures 37 inches, and its steel pivots rest in bell-metal bearings. There are 4 eyepieces which, according to a preliminary determination, magnify 80, 110, 150 and 210 times. Even the highest magnification gives perfectly sharp images. Councillor \textsc{Gauss} usually works at that magnification, although it calls for somewhat more careful adjustment in altitude because the field of view is then only 9 minutes across. The reticle consists of spider threads, seven of them vertical. Stars on the equator traverse the intervening spaces in 10 seconds. Once one has acquired the skill that observing at these short intervals requires, it is extremely agreeable to be able to observe so many thread crossings in a short time. With this greater number of threads, it also matters less if a crossing at one or several threads goes unobserved, whether by design or by chance. From several hundred observations he made of stars that are not too far from the equator, Councillor \textsc{Gauss} finds the so-called probable error of the observed crossing at one thread to be $=0^{\mathrm s}095$. On the level used to bring the axis to the horizontal, a displacement of one Paris inch corresponds to an inclination of 20 seconds. To set the telescope for any specified declination, one uses a small circle fixed to the telescope itself near the eyepiece end, divided directly into quarter degrees and by the vernier into minutes. An alidade carrying a small level rotates on it, and the index on the alidade immediately gives the declination. This arrangement has the great advantage that the instrument can be directed to any altitude with equal ease even in its reversed position. A second, similar circle to be fitted on the other side, which the instrument maker will supply later, will offer the further advantage of allowing two stars that pass through the meridian plane in quick succession to be observed without loss of time. ---
The optical perfection of the telescope is admirable. A few examples we give here from the observations of Councillor \textsc{Gauss} may show what one can make out with good, though hardly exceptionally \SourceLeaf{428}sharp eyesight, in our climate when the air is favorable but by no means exceptionally so, with due attention but without painstaking effort. With \emph{full illumination of the threads}, one of Saturn’s satellites, and often two, can readily be distinguished and observed. The same holds for the companion of the Pole Star. The little star now only 5 minutes from the pole will still bear sufficient illumination to be observed well. It can already be made out and observed at twilight, when the threads still need no illumination at all. In daylight, stars down to the 3rd magnitude can generally be observed with ease if they are neither too low nor too close to the Sun. When the air is favorable and a few precautions are taken, stars of the fourth and even the fifth magnitude can still be made out quite well in daylight. For example, Councillor \textsc{Gauss} observed the fifth-magnitude stars $\gamma'$ in Ursa Minor, $e$ in Ursa Major, $1$ in Draco according to Hevel, and $g$ in Ursa Major (Alcor) in broad daylight, the last even at noon. The difficulty of observing such faint stars in daylight really lies only in the fact that the nature of the instrument does not allow it to be set quite precisely. Where \emph{this} difficulty disappears, so that the eye need only fix on a very small part of the field of view, stars of the 5th magnitude and even fainter ones are easily made out. This is the case with double stars. For example, the companion of $\zeta$ in Ursa Major, to which the star catalogs assign only the sixth magnitude, is beautifully visible in full midday light. The still fainter companion of $\alpha$ in Hercules can likewise be made out very well at present, when it culminates at 4 in the afternoon. The many advantages of being able to observe even fainter stars well in daylight will be evident to those familiar with the subject. We mention only one, a very important one. In adjusting the instrument by means of circumpolar stars, an adjustment known to be the only one worthy of such an instrument and to be made as often as possible, one is no longer limited to the Pole Star, one of whose culminations often falls at an inconvenient hour of the night or whose observation is thwarted by unfavorable circumstances.

\par\indent The device for lifting the instrument out and reversing it also deserves special mention. The telescope is lifted out of its bearings by a rope that passes over a pulley above and leads down to a winch fixed low on the wall. The telescope is then turned around while freely suspended in the air, without the slightest jolt, and set back in its bearings. The whole operation takes only a few minutes. The pulley above the telescope the ceiling of the room is slid aside with a pull when the work is done, so that it does not get in the way when stars at the zenith are observed. This entire fine apparatus was made very practical and sturdy by our local inspector, Mr. \textsc{Rumpf}.

\par\indent We also give as examples some of the first observations made with this instrument. We note only that the declinations could not be observed on this occasion because \textsc{Repsold}’s meridian circle had been taken down during this month.

\Needspace{27\baselineskip}
\begin{center}
\emph{Observations of Saturn.}

\small
\setlength{\tabcolsep}{5pt}
\begin{tabular}{l|rrr|rrr}
\multicolumn{1}{c|}{1819}  &  \multicolumn{3}{c|}{Mean time}  &  \multicolumn{3}{c}{Apparent right asc.}  \\
\hline
Sept.\ 4  &  $13^{\mathrm u}$  &  $5^{\mathrm m}$  &  $43^{\mathrm s}9$  &  $359^\circ$  &  $40'$  &  $2''7$  \\
5  &  $13$  &  $1$  &  $22.0$  &   &  $36$  &  $2.7$  \\
6  &  $12$  &  $57$  &  $10.7$  &   &  $32$  &  $9.6$  \\
7  &  $12$  &  $52$  &  $58.8$  &   &  $28$  &  $10.3$  \\
8  &  $12$  &  $48$  &  $46.1$  &   &  $24$  &  $4.0$  \\
9  &  $12$  &  $44$  &  $34.4$  &   &  $20$  &  $0.4$  \\
10  &  $12$  &  $40$  &  $22.2$  &   &  $15$  &  $55.6$  \\
11  &  $12$  &  $36$  &  $10.1$  &   &  $11$  &  $51.7$  \\
12  &  $12$  &  $31$  &  $57.5$  &   &  $7$  &  $41.2$  \\
\hline
 \multicolumn{1}{c|}{1819}  &  \multicolumn{3}{c|}{Mean time}  &  \multicolumn{3}{c}{Apparent right asc.}  \\
\hline
 Sept.\ 13  &  $12^{\mathrm u}$  &  $27^{\mathrm m}$  &  $45^{\mathrm s}0$  &  $359^\circ$  &  $3'$  &  $30''4$  \\
 14  &  $12$  &  $23$  &  $32.3$  &  $358$  &  $59$  &  $18.1$  \\
 15  &  $12$  &  $19$  &  $19.8$  &   &  $55$  &  $9.0$  \\
 16  &  $12$  &  $15$  &  $7.0$  &   &  $50$  &  $54.4$  \\
 19  &  $12$  &  $2$  &  $28.6$  &   &  $38$  &  $10.6$  \\
 21  &  $11$  &  $54$  &  $3.0$  &   &  $29$  &  $42.9$  \\
 24  &  $11$  &  $41$  &  $24.2$  &   &  $15$  &  $56.4$  \\
 28  &  $11$  &  $24$  &  $33.4$  &   &  $0$  &  $5.1$  \\
 Oct.\ 1  &  $11$  &  $11$  &  $55.9$  &  $357$  &  $47$  &  $36.3$  \\
\end{tabular}
\end{center}

\Needspace{22\baselineskip}
\SourceLeaf{429}
\emph{Observations of Vesta.}

\begin{center}
\begin{tabular}{l|rrr|rrr}
\multicolumn{1}{c|}{1819}  &  \multicolumn{3}{c|}{Mean time}  &  \multicolumn{3}{c}{Apparent right asc.}  \\
\hline
Sept.\ 7  &  $13^{\mathrm u}$  &  $32^{\mathrm m}$  &  $52^{\mathrm s}0$  &  $9^\circ$  &  $28'$  &  $7''2$  \\
9  &   &  $23$  &  $32.6$  &   &  $6$  &  $10.5$  \\
10  &   &  $18$  &  $51.1$  &  $8$  &  $54$  &  $44.7$  \\
11  &   &  $14$  &  $8.6$  &   &  $43$  &  $1.8$  \\
12  &   &  $9$  &  $25.1$  &   &  $31$  &  $6.9$  \\
15  &  $12$  &  $55$  &  $18.7$  &  $7$  &  $53$  &  $49.8$  \\
\hline
 \multicolumn{1}{c|}{1819}  &  \multicolumn{3}{c|}{Mean time}  &  \multicolumn{3}{c}{Apparent right asc.}  \\
\hline
 Sept.\ 16  &  $12^{\mathrm u}$  &  $50^{\mathrm m}$  &  $21^{\mathrm s}4$  &  $7^\circ$  &  $40'$  &  $58''5$  \\
 19  &   &  $35$  &  $55.5$  &   &  $1$  &  $18.0$  \\
 21  &   &  $26$  &  $15.9$  &  $6$  &  $34$  &  $15.9$  \\
 24  &   &  $11$  &  $43.4$  &  $5$  &  $52$  &  $58.0$  \\
 28  &  $11$  &  $52$  &  $18.1$  &  $4$  &  $57$  &  $24.3$  \\
 Oct.\ 1  &   &  $37$  &  $45.3$  &   &  $16$  &  $1.2$  \\
\end{tabular}
\end{center}

\Needspace{27\baselineskip}
\par\indent Since these observations have yet to be compared with tables or elements, we also give the observations of a fixed star, together with their reduction to the mean position for the beginning of the year 1820.


\begin{center}
\emph{Right ascension of 85 Pegasi.}
\end{center}

\begin{center}
\begin{tabular}{l|rrr|rrr}
\multicolumn{1}{c|}{1819}  &  \multicolumn{3}{c|}{Observed}  &  \multicolumn{3}{c}{Mean for 1820}  \\
\hline
Sept.\ 4  &  $358^\circ$  &  $11'$  &  $40''8$  &  $358^\circ$  &  $11'$  &  $43''0$  \\
5  &   &   &  $41.4$  &   &   &  $43.4$  \\
6  &   &   &  $40.2$  &   &   &  $42.0$  \\
7  &   &   &  $43.9$  &   &   &  $45.5$  \\
8  &   &   &  $42.6$  &   &   &  $44.0$  \\
9  &   &   &  $42.9$  &   &   &  $44.1$  \\
10  &   &   &  $43.0$  &   &   &  $44.1$  \\
\hline
 \multicolumn{1}{c|}{1819}  &  \multicolumn{3}{c|}{Observed}  &  \multicolumn{3}{c}{Mean for 1820}  \\
\hline
 Sept.\ 11  &  $358^\circ$  &  $11'$  &  $42''6$  &  $358^\circ$  &  $11'$  &  $43''6$  \\
 12  &   &   &  $45.4$  &   &   &  $46.2$  \\
 13  &   &   &  $41.8$  &   &   &  $42.4$  \\
 14  &   &   &  $44.7$  &   &   &  $45.1$  \\
 15  &   &   &  $44.8$  &   &   &  $45.1$  \\
 16  &   &   &  $43.0$  &   &   &  $43.2$  \\
\multicolumn{4}{r|}{Mean . .} &  $358$  &  $11$  &  $43.97$  \\
\end{tabular}
\end{center}

\par\indent This star is notable for its large proper motion. Comparing its right ascension as derived from these observations with \textsc{Bradley}’s for 1755 gives a proper motion in 64.7 years of $+54''65$, and hence in one year of $+0''845$.
